%=============================================================
% Capitolo 1 — Aritmetica in Virgola Mobile
%=============================================================
\chapter{Aritmetica in Virgola Mobile}
\label{ch:fp}

\section{Introduzione: il problema della rappresentazione}

Nei calcolatori moderni i numeri reali vengono rappresentati con un numero
\emph{finito} di cifre binarie. Questo comporta due conseguenze fondamentali:
\begin{itemize}
  \item non tutti i numeri reali sono rappresentabili esattamente;
  \item le operazioni aritmetiche introducono errori di arrotondamento.
\end{itemize}

Comprendere questi limiti è essenziale per valutare l'affidabilità dei risultati
numerici.

%------------------------------------------------------------
\section{Sistema posizionale e notazione scientifica}
%------------------------------------------------------------

Un numero reale $x$ viene espresso in \emph{notazione scientifica normalizzata} in
base $b$ come:
\begin{equation}
  x = \pm\, m \cdot b^e,
  \label{eq:scientific}
\end{equation}
dove:
\begin{itemize}
  \item $b \geq 2$ è la \emph{base} del sistema (tipicamente $b = 2$ per i calcolatori);
  \item $m$ è la \emph{mantissa} (o significando), con $1 \leq m < b$;
  \item $e$ è l'\emph{esponente} (intero con segno).
\end{itemize}

\begin{example}[Notazione scientifica in base 10]
Il numero $\pi \approx 3{,}141592\ldots$ si scrive come
\[
  3{,}141592 \times 10^{0}.
\]
Con 5 cifre disponibili per la mantissa, la rappresentazione è
$3{,}14159 \times 10^{0}$ (troncamento) oppure $3{,}14159 \times 10^{0}$ (arrotondamento).
\end{example}

%------------------------------------------------------------
\section{Lo standard IEEE 754}
%------------------------------------------------------------

Lo standard \textbf{IEEE 754} (1985, aggiornato nel 2008 e 2019) definisce il formato
dei numeri in virgola mobile usati in praticamente tutti i processori moderni.

\subsection{Formato della parola binaria}

Un numero floating-point è memorizzato in una parola di $w$ bit suddivisa in tre campi:

\begin{center}
\begin{tabular}{|c|c|c|}
\hline
\textbf{Segno} (1 bit) & \textbf{Esponente} ($e$ bit) & \textbf{Mantissa} ($m$ bit) \\
\hline
\end{tabular}
\end{center}

\medskip
I due formati principali sono riassunti in \cref{tab:ieee754}.

\begin{table}[htbp]
  \centering
  \caption{Formati IEEE 754 a precisione singola e doppia.}
  \label{tab:ieee754}
  \begin{tabular}{lcccc}
    \toprule
    \textbf{Formato} & \textbf{Bit totali} & \textbf{Bit segno} &
    \textbf{Bit esponente} & \textbf{Bit mantissa} \\
    \midrule
    Precisione singola & 32 & 1 & 8  & 23 \\
    Precisione doppia  & 64 & 1 & 11 & 52 \\
    \bottomrule
  \end{tabular}
\end{table}

\subsection{Il bit nascosto}

Nella notazione normalizzata la prima cifra della mantissa è sempre 1 (in base 2),
quindi non viene memorizzata esplicitamente: si guadagna così un bit di precisione
\emph{gratuitamente}. Questo bit è detto \textbf{bit nascosto} (o \emph{implicit
leading bit}).

La mantissa effettiva è quindi $1.f_1f_2\ldots f_m$ dove $f_i \in \{0,1\}$.

\subsection{Configurazioni speciali}

Con $e_\text{bit}$ bit per l'esponente si hanno $2^{e_\text{bit}}$ configurazioni
possibili. Due di queste sono riservate a valori speciali:
\begin{itemize}
  \item \texttt{+Inf}, \texttt{-Inf} — overflow;
  \item \texttt{NaN} (Not a Number) — operazione non definita (es.\ $0/0$).
\end{itemize}
Le rimanenti $2^{e_\text{bit}} - 2$ codificano gli esponenti normali. Ad esempio in
precisione singola si hanno $256 - 2 = 254$ esponenti utilizzabili.

\subsection{Numeri denormalizzati}

I \textbf{numeri denormalizzati} (o subnormali) permettono di rappresentare valori
molto piccoli, al di sotto del minimo normalizzato, a prezzo di una riduzione della
precisione. In questo caso il bit nascosto viene posto a 0:
\[
  x = \pm\, 0.f_1f_2\ldots f_m \times b^{e_\min}.
\]

%------------------------------------------------------------
\section{Insieme dei numeri di macchina}
%------------------------------------------------------------

\begin{definition}[Insieme dei numeri di macchina]
L'insieme $\F \subset \R$ dei \emph{numeri di macchina} in base $b$, con $t$ cifre di
mantissa ed esponente nell'intervallo $[e_\min, e_\max]$, è:
\begin{equation}
  \F(b, t, e_\min, e_\max) = \bigl\{ x \in \R \;:\; x = \pm\, b^e \sum_{i=1}^{t} d_i\, b^{-i},\;
  e \in [e_\min, e_\max],\; d_i \in \{0,\ldots,b{-}1\},\; d_1 \neq 0 \bigr\} \cup \{0\}.
  \label{eq:macchina}
\end{equation}
\end{definition}

\begin{remark}
$\F$ è un insieme \emph{finito} e \emph{discreto}: i numeri di macchina non sono
distribuiti uniformemente sulla retta reale, ma più densamente vicino allo zero e più
rade man mano che ci si allontana.
\end{remark}

%------------------------------------------------------------
\section{Troncamento e arrotondamento}
%------------------------------------------------------------

Sia $x \in \R$ un numero non rappresentabile esattamente in $\F$. L'operazione
$\fl : \R \to \F$ che associa a $x$ il numero di macchina più vicino si chiama
\textbf{arrotondamento}.

\subsection{Troncamento}

Il \emph{troncamento} della mantissa a $t$ cifre è:
\begin{equation}
  \fl_\text{tronc}(x) = \text{sign}(x) \cdot b^e \cdot \sum_{i=1}^{t} d_i\, b^{-i}.
\end{equation}

\subsection{Arrotondamento al più vicino}

L'\emph{arrotondamento al più vicino} (default IEEE) sceglie il numero di macchina più
prossimo a $x$:
\begin{equation}
  \fl_\text{arr}(x) = \arg\min_{y \in \F} |x - y|.
\end{equation}

In caso di equidistanza si applica la regola \emph{round half to even} (arrotondamento
bancario).

%------------------------------------------------------------
\section{Errore di arrotondamento ed epsilon di macchina}
%------------------------------------------------------------

\begin{definition}[Errore relativo di rappresentazione]
L'errore relativo commesso nell'approssimare $x \neq 0$ con $\fl(x)$ è:
\begin{equation}
  \delta = \frac{\fl(x) - x}{x}.
  \label{eq:errore_rel}
\end{equation}
\end{definition}

\begin{definition}[Epsilon di macchina]
L'\emph{epsilon di macchina} $\eps_M$ è il più piccolo numero positivo tale che
\begin{equation}
  \fl(1 + \eps_M) > 1.
  \label{eq:eps}
\end{equation}
Equivalentemente, $\eps_M$ è la distanza relativa tra 1 e il successivo numero di
macchina.
\end{definition}

Per il formato IEEE 754 si ha:
\begin{align}
  \eps_M^{(32)} &= 2^{-23} \approx 1.19 \times 10^{-7} \quad \text{(singola precisione)}, \\
  \eps_M^{(64)} &= 2^{-52} \approx 2.22 \times 10^{-16} \quad \text{(doppia precisione)}.
\end{align}

\begin{theorem}[Limitazione dell'errore relativo]
\label{thm:errore_rel}
Per ogni $x \in \R$ con $|x|$ nel range dei numeri normalizzati vale:
\begin{equation}
  |\delta| = \left|\frac{\fl(x) - x}{x}\right| \leq \frac{1}{2}\, b^{1-t} = \frac{\eps_M}{2}.
  \label{eq:bound_errore}
\end{equation}
Equivalentemente, $\fl(x) = x(1 + \delta)$ con $|\delta| \leq \eps_M/2$.
\end{theorem}

\begin{remark}[Underflow]
Se un numero $x \neq 0$ è troppo piccolo in valore assoluto per essere rappresentato
come normalizzato (e anche come denormalizzato), si ha \emph{underflow}: il numero
viene arrotondato a zero. Questo non è la stessa cosa di $x = 0$.
\end{remark}

%------------------------------------------------------------
\section{Algoritmo di Horner}
\label{sec:horner}
%------------------------------------------------------------

La valutazione di un polinomio
\begin{equation}
  p(x) = a_n x^n + a_{n-1} x^{n-1} + \cdots + a_1 x + a_0
  \label{eq:polinomio}
\end{equation}
richiede, con la formula diretta, $n$ addizioni e $\frac{n(n+1)}{2}$ moltiplicazioni.

L'\textbf{algoritmo di Horner} (o schema di Ruffini–Horner) riduce il costo a $n$
addizioni e $n$ moltiplicazioni, riscrivendo il polinomio in forma annidiata:
\begin{equation}
  p(x) = a_0 + x\bigl(a_1 + x\bigl(a_2 + \cdots + x(a_{n-1} + x\, a_n)\cdots\bigr)\bigr).
  \label{eq:horner}
\end{equation}

\begin{algorithm}[htbp]
  \caption{Algoritmo di Horner per la valutazione di polinomi}
  \label{alg:horner}
  \begin{algorithmic}[1]
    \Require Coefficienti $a_0, a_1, \ldots, a_n$; punto di valutazione $x$
    \Ensure $p \approx p(x)$
    \State $p \gets a_n$
    \For{$k = n-1$ \textbf{downto} $0$}
      \State $p \gets p \cdot x + a_k$
    \EndFor
    \State \Return $p$
  \end{algorithmic}
\end{algorithm}

Il costo computazionale dell'Algoritmo~\ref{alg:horner} è $\mathcal{O}(n)$.

%------------------------------------------------------------
\section{Propagazione degli errori}
\label{sec:propagazione}
%------------------------------------------------------------

Sia $f : \R^n \to \R$ una funzione differenziabile. Se i dati $x_1,\ldots,x_n$ sono
affetti da errori assoluti $\Delta x_1,\ldots,\Delta x_n$, l'errore assoluto sul
risultato è approssimato dalla formula di propagazione:
\begin{equation}
  \Delta f \approx \sum_{i=1}^{n} \left|\frac{\partial f}{\partial x_i}\right| \Delta x_i.
  \label{eq:propagazione}
\end{equation}

In termini di errori relativi (con $\Delta x_i = |\eps_i| |x_i|$):
\begin{equation}
  \frac{\Delta f}{|f|} \approx \sum_{i=1}^{n} \left|\frac{\partial f}{\partial x_i}\right|
  \frac{|x_i|}{|f|} |\eps_i|.
  \label{eq:propagazione_rel}
\end{equation}

\begin{example}[Sottrazione catastrofica]
La sottrazione di due numeri quasi uguali $a \approx b$ amplifica enormemente l'errore
relativo. Se $a = 1{,}0001$ e $b = 1{,}0000$, entrambi con errore relativo $10^{-4}$,
l'errore relativo sulla differenza $a - b = 0{,}0001$ può raggiungere il 100\%.
Questo fenomeno è noto come \emph{cancellazione numerica}.
\end{example}

%------------------------------------------------------------
\section{Aritmetica in virgola mobile IEEE}
\label{sec:aritmetica_ieee}
%------------------------------------------------------------

Lo standard IEEE 754 garantisce che le operazioni elementari $+,-,\times,\div$ (e la
radice quadrata) siano eseguite come se il calcolo fosse fatto con precisione infinita
e poi arrotondato al più vicino numero di macchina:
\begin{equation}
  \fl(x \circ y) = \fl(x \circ_\text{esatta} y) = (x \circ y)(1 + \delta),
  \quad |\delta| \leq \frac{\eps_M}{2},
  \label{eq:ieee_op}
\end{equation}
dove $\circ \in \{+,-,\times,\div\}$.

\begin{remark}[Registri estesi]
Alcuni processori usano registri interni a 80 bit per i calcoli intermedi, riducendo
l'accumulo degli errori di arrotondamento.
\end{remark}
